\section{Transferring a finite measure and a bounded metric}
\label{sec:transfer}

The next proposition isolates the induction step. It does not assume
the existence of a section on the base: it constructs base data to
which a lower-dimensional nonvanishing theorem can be applied.
We retain two models of the total space because they do different jobs.
On a normal equidimensional model every vertical prime maps onto a
base prime, so divisor orders can be normalized there. A smooth
resolution carries canonical forms and fiber integrals, but need not
remain equidimensional. Such diagrams arise from flattening a resolved
rational map \cite{RaynaudGruson1971}.

A prime is called \emph{horizontal} if it dominates the base, and
\emph{vertical} otherwise. For a morphism \(f:Y\to S\) of smooth varieties
we write \(K_{Y/S}=K_Y-f^*K_S\). Its restriction to a smooth fiber is the
canonical line of that fiber.

\begin{proposition}[Two-metric transfer]\label{prop:transfer}
Let \(X,S,Y\) be smooth connected projective varieties over \(\C\), and
let \(\bar Y\) be a normal projective variety. Suppose
\(r:Y\to\bar Y\) and \(\bar\pi:\bar Y\to X\) are projective birational
morphisms, and \(\bar f:\bar Y\to S\) is a surjective equidimensional
morphism with connected general fiber. Put
\(\pi=\bar\pi r\) and \(f=\bar f r\), as in Figure~\ref{fig:models}.

Let \(D\) be an effective integral divisor on \(X\), and let
\(J=-K_X+D\) have a semipositive singular metric \(h_J\) with
\[
 \int_X |s_D|_{h_J}^2<\infty.
\]
Let \(L\) be a line bundle on \(X\) with a semipositive metric \(h_L\)
of locally bounded weights. Suppose that for an integer \(k>0\)
there is a nonzero rational section \(\rho\) of \(k\bar\pi^*L\)
with no horizontal poles, and that
\begin{equation}\label{eq:rank}
 \dim_{\C(S)}
 H^0\!\left(Y_\eta,
  \left(K_{Y/S}+\pi^*(J+jkL)\right)|_{Y_\eta}\right)=1
 \qquad(j=0,1,2,\ldots).
\end{equation}
Here \(\eta\) is the generic point of \(S\) and
\(Y_\eta=Y\times_S\operatorname{Spec}\C(S)\).

After replacing \(k,\rho\) by \(dk,\rho^d\) for some integer \(d>0\),
there are a line bundle \(A\) on \(S\), an effective integral divisor
\(D_S\), and a nonzero regular section
\[
 e\in H^0(\bar Y,k\bar\pi^*L-\bar f^*A)
\]
with these properties:
\begin{enumerate}[label=\textup{(\roman*)}]
\item \(A\) has a semipositive metric with locally bounded weights.
\item \(-K_S+D_S\) has a semipositive singular metric \(h_S\) with
      \(\int_S|s_{D_S}|_{h_S}^2<\infty\).
\item \(D_S\) is supported on base primes \(T\) for which every prime
      component of \(\bar f^*T\) is \(\bar\pi\)-exceptional or lies in
      \(\Supp(\bar\pi^*D)\).
\item If \(a>0,b\ge0\) are integers and \(H^0(S,aA+bD_S)\ne0\), then
      \(H^0(X,akL+cD)\ne0\) for some integer \(c\ge0\).
\end{enumerate}
In these conclusions \(k\) denotes the enlarged integer.
\end{proposition}

\begin{figure}[ht]
\centering
\begin{tikzcd}[row sep=large,column sep=large]
Y \arrow[r,"r"] \arrow[dr,"f"'] \arrow[rr,bend left=22,"\pi"]
  & \bar Y \arrow[r,"\bar\pi"] \arrow[d,"\bar f"]
  & X\\
  & S &
\end{tikzcd}
\caption{The model diagram in Proposition~\ref{prop:transfer}.
The normal space \(\bar Y\) is equidimensional over \(S\); the smooth
space \(Y\) is used for integration. Both upper maps are birational.
The diagram does not assert that \(f\) is equidimensional.}
\label{fig:models}
\end{figure}

\begin{proof}
We first normalize \(\rho\) so that a divided section is regular and
nonzero on at least one component over each base prime. We then use
its adjoint integrals twice: the zeroth integral gives the volume
metric, and the high moments give the bounded metric. The support
of the first correction will make the final lifting possible.

If \(S\) is a point, every prime of \(\bar Y\) is horizontal.
Thus \(\rho\) is regular by normality and descends as a nonzero
section of \(kL\) on the smooth variety \(X\).
Take \(A\) trivial and \(D_S=0\); the metric assertions on a point
and the lifting assertion follow. We assume \(\dim S>0\) below.

\medskip\noindent
\emph{Divisorial normalization.}
For a prime divisor \(T\subset S\), define
\begin{equation}\label{eq:minima}
 c_T=\min_{\bar f(Q)=T}
       \frac{\ord_Q(\rho)}{\ord_Q(\bar f^*T)},
\end{equation}
where \(Q\) ranges over the prime divisors of \(\bar Y\) mapping onto
\(T\). These are finitely many components for each \(T\), and the
denominators are positive integers. Only finitely many \(c_T\) are
nonzero. Replace \(\rho\) by a power and \(k\) by the corresponding
multiple to make every \(c_T\) integral, and set
\(A=\OO_S(\sum_Tc_TT)\). The rank condition \eqref{eq:rank} persists:
the new adjoint twists form a subsequence of the old ones and still
include \(j=0\).

Divide \(\rho\) by the pullback of the tautological rational section
of \(A\), obtaining a rational section \(e\) of
\(k\bar\pi^*L-\bar f^*A\).
Its order is nonnegative at horizontal primes by hypothesis and at
vertical primes by \eqref{eq:minima}. Equidimensionality ensures that
every vertical prime maps onto a base prime. Thus these are all prime
divisors of \(\bar Y\), and normality makes \(e\) regular.
For each base prime \(T\), a component attaining the minimum in
\eqref{eq:minima} has order zero in \(e\).
We also write \(e\) for its pullback to \(Y\).

\medskip\noindent
\emph{The fiber integrals and their curvature.}
Pullback of the \(J\)-valued top form includes the Jacobian:
\[
 \eta_Y=\pi^*s_D\in H^0(Y,K_Y+\pi^*J),\qquad
 \mu_Y=|\eta_Y|_{\pi^*h_J}^2.
\]
Change of variables gives finite total mass for \(\mu_Y\).
In particular this assertion concerns measures and does not require
a pointwise value for a singular metric on its polar set.

On a coordinate chart \(U\subset S\), take a canonical frame \(ds\)
of \(K_S\) and a frame \(v\) of \(A\). The section
\(e_v=e f^*v\) belongs to \(k\pi^*L\) over \(f^{-1}(U)\). For
\(j=0,1,2,\ldots\), set
\begin{equation}\label{eq:moments}
 u_j=(\eta_Y/f^*ds)e_v^j,\qquad
 V_j(s)=\int_{Y_s}|u_j|_{\pi^*(h_Jh_L^{jk})}^2,
 \qquad \Phi_j=-\log V_j .
\end{equation}
The section \(u_j\) belongs to \(K_{Y/S}+\pi^*(J+jkL)\).
On a smooth fiber it is a top form with coefficients in the indicated
twisting line; \eqref{eq:moments} pairs that form with its conjugate
to give a positive density. Properness, regularity of \(e_v\), and
boundedness of \(h_L\)'s weights bound \(|e_v|^2\) on inverse images
of relatively compact base charts. Since \(V_0\) is the density of
the finite measure \(f_*\mu_Y\), all \(V_j\) are finite and positive
almost everywhere on a suitable smooth-fibration open.

For clarity, changes of frames have the following effect:
\begin{equation}\label{eq:frames}
 v\mapsto gv,\quad ds\mapsto h\,ds
 \quad\Longrightarrow\quad
 \Phi_j\mapsto \Phi_j-j\log|g|^2+\log|h|^2 .
\end{equation}
This will give weights on two different lines: \(-K_S\) at \(j=0\),
and \(A\) after division by \(j\) and passage to the limit.

The relative Bergman theorem
\cite[Theorem~0.1]{BerndtssonPaun2008}
applies to the semipositive twisting metric
\(\pi^*(h_Jh_L^{jk})\): \eqref{eq:moments} supplies a nonzero
square-integrable adjoint section on almost every smooth fiber.
The theorem produces a semipositive relative Bergman metric.
On a dense base open where base change holds and \eqref{eq:rank}
has dimension one, \(u_j\) generates the adjoint section space.
For almost every such fiber \(V_j<\infty\), so its one-dimensional
holomorphic adjoint space is also its full space of square-integrable
sections: every section is a scalar multiple of \(u_j\).
Consequently the Bergman weight, in a holomorphic frame upstairs, is
\[
 \Phi_j\circ f+\log|\widehat u_j|^2
\]
almost everywhere, where \(\widehat u_j\) is the coefficient of \(u_j\).
Off the zero locus of that coefficient, its logarithm is
pluriharmonic, so subtracting it leaves a psh function upstairs.
In a local product chart, Fubini's theorem gives almost all constant
slices on which the displayed identity holds almost everywhere.
Restriction to any of these slices gives a psh representative of
\(\Phi_j\) on the base. For fixed base frames, these representatives
agree almost everywhere on overlaps, hence everywhere, and therefore patch.

The base-change open can depend on \(j\), so we now put all these
weights on one fixed open \(S^\circ\). Choose \(S^\circ\) such that
\(f\) is smooth there and every fiber has a point where both
\(\eta_Y\) and \(e\) are nonzero. A psh weight is locally bounded above;
therefore its squared frame norm has a positive local lower bound,
even when the metric is singular. On a small product chart around
such a point, this lower bound for \(h_J\), together with the
bounds for \(h_L\), gives
\begin{equation}\label{eq:common-open}
 V_j(s)\ge C_1C_2^j,\qquad \Phi_j(s)\le C+jC',
\end{equation}
locally on \(S^\circ\), with positive \(C_1,C_2\) independent of \(j\).
Removable singularities extend each psh representative across its
additional exceptional base loci inside \(S^\circ\).
The union of the exceptional measure-zero sets for the countably many
\(j\) still has measure zero; hence all moment comparisons can be used
simultaneously. Almost-everywhere upper bounds also hold for the psh
representatives by the submean inequality.

We have now obtained all the moment weights on a fixed open set.
To continue the induction, both kinds of weight must extend to the
boundary. Their estimates use different components of a singular fiber.

\medskip\noindent
\emph{The zeroth moment: finite volume with supported correction.}
Call a base prime \(T\) \emph{bad} if every component of
\(\bar f^*T\) is \(\bar\pi\)-exceptional or lies in
\(\Supp(\bar\pi^*D)\). Only finitely many primes are bad: each is the
image of a member of the finite list of exceptional and boundary
divisors on \(\bar Y\).

Take a prime \(T\) in the complement of \(S^\circ\), and a component
\(Q\) above it on \(\bar Y\). At general points of its strict transform
in \(Y\), there are coordinates
\[
 s_1=t^\ell,\qquad s_i=x_i\quad(2\le i\le\dim S),
 \qquad \ell=\ord_Q(\bar f^*T)>0 .
\]
Indeed, the map on the divisors is generically smooth, and a root of
a nonvanishing unit absorbs the unit in the first coordinate.
Write the coefficient of \(\eta_Y\) as \(t^q\) times a unit at such
a point, with \(q\ge0\). Division by \(ds\) contributes
\(t^{-(\ell-1)}\), so the relative top-form coefficient has magnitude
comparable to \(|t|^{q-\ell+1}\). The metric lower bound and integration
over a fixed patch in the free fiber coordinates give
\begin{equation}\label{eq:zero-bound}
 V_0(s)\ge C|s_1|^{2(q-\ell+1)/\ell}
\end{equation}
almost everywhere near a general point of \(T\).

If \(T\) is not bad, choose \(Q\) nonexceptional over \(X\) and
outside \(\Supp(\bar\pi^*D)\). Then the birational Jacobian and \(s_D\)
are generically nonzero on \(Q\), so \(q=0\). Thus \(-\log V_0\)
is locally bounded above there. At the bad primes choose sufficiently
large nonnegative integral coefficients for \(D_S\). Equation
\eqref{eq:zero-bound} makes the weights
\begin{equation}\label{eq:base-volume-weight}
 \Phi_0+\log|\widehat s_{D_S}|^2
\end{equation}
locally bounded above at general points of every boundary prime.
Here the hat denotes the coefficient in a local frame of \(\OO_S(D_S)\).
These weights are psh off the boundary, extend across its general
divisorial points, and then extend across the remaining analytic set
of codimension at least two. The frame rule \eqref{eq:frames} gives a
metric on \(-K_S+D_S\).

Pairing \eqref{eq:base-volume-weight} with \(s_{D_S}\) cancels the
added factor. The resulting adjoint measure equals \(f_*\mu_Y\)
almost everywhere, so has finite total mass. We have obtained the
volume metric and the required support of its correction.

\medskip\noindent
\emph{High moments: bounded weights on \(A\).}
For \(l=1,2,\ldots\), put on \(S^\circ\)
\begin{equation}\label{eq:envelope}
 \Psi=\lim_{l\to\infty}
      \left(\sup_{j\ge l}\frac{\Phi_j}{j}\right)^* ,
\end{equation}
where the star denotes upper-semicontinuous regularization in fixed
local frames. Equation \eqref{eq:common-open} gives uniform local
upper bounds. The regularized suprema are therefore psh and decrease.

Their limit does not collapse to \(-\infty\).
On any relatively compact base chart \(U\), including charts meeting
the boundary, properness and bounded weights of \(h_L\) give a single
constant \(C_U>0\) with
\begin{equation}\label{eq:uniform-lower}
 |e_v|_{\pi^*h_L^k}^2\le C_U,\qquad V_j\le C_U^jV_0 .
\end{equation}
The first bound holds on the whole inverse image of \(U\):
take a finite cover of its compact closure by line-bundle frames.
For almost every \(s\in U\cap S^\circ\), \(V_0(s)\) is finite and
positive. Dividing \eqref{eq:uniform-lower} by \(j\) after taking
negative logarithms gives
\(\Psi\ge-\log C_U\).
Thus \(\Psi\) is psh. Its frame transformation is exact, although
that of each normalized moment has a vanishing error. On a relatively
compact overlap, \eqref{eq:frames} changes \(\Phi_j/j\) by
\(-\log|g|^2+(\log|h|^2)/j\). The last term has absolute value
at most \(C/l\) for every \(j\ge l\). The two tail suprema,
and then their upper-semicontinuous regularizations, therefore differ
from one another plus \(-\log|g|^2\) by at most \(C/l\).
Taking the decreasing limit gives precisely
\(\Psi\mapsto\Psi-\log|g|^2\), the weight rule for \(A\).

To extend \(\Psi\), now choose a component \(Q\) attaining the
minimum in \eqref{eq:minima}. This need not be the component used
for the zeroth moment. Its defining property is that \(e\) is nonzero
generically on \(Q\). The local norm of \(e_v\) is therefore bounded
below on a suitable patch. The same power-map calculation as in
\eqref{eq:zero-bound}, with the factor \(e_v^j\), gives
\begin{equation}\label{eq:high-bound}
 V_j(s)\ge C_1C_2^j|s_1|^{2b_T}
\end{equation}
for a nonnegative integer \(b_T\) independent of \(j\).
For example, increasing a nonnegative integer above
\((q-\ell+1)/\ell\) suffices after shrinking the coordinate chart.
The absence of zeros of \(e\) on the chosen component is what prevents
an additional exponent proportional to \(j\).

Shrink the chart so that \(|s_1|<1\), and decrease \(C_1\) if needed
so that \(C_1\le1\). Put
\[
 B(s)=-\log C_1-2b_T\log|s_1|\ge0.
\]
For \(j\ge l\), the negative logarithm of \eqref{eq:high-bound} gives
\[
 \frac{\Phi_j(s)}{j}\le-\log C_2+\frac{B(s)}{j}
                  \le-\log C_2+\frac{B(s)}{l}.
\]
The right side is continuous off \(T\), so it also bounds the
upper-semicontinuous regularization of the tail supremum there.
Letting \(l\to\infty\) in \eqref{eq:envelope} gives
\(\Psi\le-\log C_2\) throughout the punctured chart.
Hence \(\Psi\) is locally bounded above near general points of
each boundary prime, and extends psh across them and then across
codimension two. The lower bound in \eqref{eq:uniform-lower} was
uniform on the whole base chart, so the extension retains it.
The extended psh weight is locally bounded above as well.
Thus the metric on \(A\) has locally bounded weights on both sides.

\medskip\noindent
\emph{Lifting the section.}
We have the two metrics needed on the base. It remains to verify
that the divisor used for finite volume creates only permitted poles
on \(X\). Regard a nonzero section of \(aA+bD_S\) as a rational section
of \(aA\) with poles only on \(D_S\). Pull it to the normal
equidimensional model \(\bar Y\) and multiply by \(e^a\).
The resulting nonzero rational section of \(ak\bar\pi^*L\) has poles
only above bad primes. By their definition every such prime divisor
of \(\bar Y\) is exceptional over \(X\) or lies above \(D\).

At every prime of \(X\) outside \(D\), the strict transform in
\(\bar Y\) therefore has no pole. The rational section downstairs
extends there; exceptional divisors give no additional obstruction
on the smooth, hence normal, variety \(X\).
Its remaining finitely many pole orders are cleared by multiplication
by \(s_D^c\) for a sufficiently large integer \(c\ge0\).
This gives a nonzero section of \(akL+cD\), as required.
\end{proof}
